\paper{非数学 B 类}

\begin{problems}
\item[一、]（本题 30 分，每小题 6 分）填空题：
\begin{enumerate}
  \item 极限 $\displaystyle \lim_{x\to 0}\frac{\ln(1+x^2)-x^2}{\sin^4 x}=\underline{\hspace{3.5em}}$．
  \item $\displaystyle \int_{-2}^{2}(x^3\cos^5 x+\sqrt{4-x^2})\dd x=\underline{\hspace{3.5em}}$．
  \item 设 $y=y(x)$ 由方程 $\displaystyle \ee^{3y}+\int_0^{x+y}\cos t^2\dd t=1$ 确定．则 $\displaystyle \left.\frac{\dd^2 y}{\dd x^2}\right|_{(0,0)}=\underline{\hspace{3.5em}}$．
  \item 设 $z=f(xy,\ee^{x+y})$，其中 $f$ 具有二阶连续偏导数．则
  \[
    z_{xy}=\underline{\hspace{10em}}.
  \]
  \item 设 $D:\ x^2+y^2\le r^2$，其中 $r>0$．则 $\displaystyle \lim_{r\to 0^+}\frac{\displaystyle \iint_D (\ee^{x^2+y^2}-1)\dd x\dd y}{r^4}=\underline{\hspace{3.5em}}$．
\end{enumerate}

\item[二、]（本题 14 分）用 Lagrange 乘子法求函数 $f(x,y,z)=x+2y+3z$ 在平面 $x-y+z=1$ 和圆柱面 $x^2+y^2=1$ 的交线上的最大值．

\item[三、]（本题 14 分）求函数 $\displaystyle f(x)=\frac{1}{(x-1)(x+3)}$ 在 $x_0=2$ 处的 Taylor 级数，并确定它的收敛域．

\item[四、]（本题 14 分）求微分方程 $(x^3-y^2)\dd x+(x^2y+xy)\dd y=0$ 的通解．

\item[五、]（本题 14 分）已知 $\displaystyle \int_0^{+\infty}\frac{\sin x}{x}\dd x=\frac{\pi}{2}$．
\begin{enumerate}
  \item 计算 $\displaystyle I_1=\int_0^{+\infty}\left(\frac{\sin x}{x}\right)^2\dd x$．
  \item 计算 $\displaystyle I_2=\int_0^{+\infty}\left(\frac{\sin x}{x}\right)^3\dd x$．
\end{enumerate}

\item[六、]（本题 14 分）设 $f(x)$ 是 $(-\infty,+\infty)$ 上具有连续导数的非负函数，且存在 $M>0$ 使得对任意的 $x,y\in(-\infty,+\infty)$，有 $|f'(x)-f'(y)|\le M|x-y|$．证明：对于任意实数 $x$，恒有 $(f'(x))^2\le 2Mf(x)$．
\end{problems}
